\section{はじめに}\label{sec:intro}

コヒーレント・イジングマシン（CIM）は，光パラメトリック発振器のネットワークでイジング模型の基底状態を探す物理的な計算機である\cite{wang2013,inagaki2016,mcmahon2016}．MAX-CUT のような組合せ最適化問題を高速に近似的に解けるが，大規模な問題では出力が最適解から離れることが多く，実用上は CIM の出力を初期解として，焼きなまし法（SA）\cite{kirkpatrick1983}やメメティックアルゴリズム\cite{wu2012}などの後段の探索で磨くハイブリッドな使い方が考えられる．

このとき CIM は多数の候補を安価に生成できるので，「どの候補を後段に渡すか」が問題になる．自然な方法はカット値の高い候補を選ぶことである．しかし，同じカット値の候補でも，後段の探索で最終的に到達できる品質が同じとは限らない．候補が探索空間のどこにあるか（最適解への近さや，そこへ向かう経路の有無）によって，後段の探索の行き着く先が変わるはずである．

本論文はこの「種の価値」の差を実測し，それを安価に予測できるか，予測に基づく選別が実用上役に立つかを調べる．具体的には次の 3 つを問う．(C1) 同じカット値の CIM 解で後段の価値に差があるか．(C2) その差をカット値以外の安価な指標で予測できるか．(C3) 指標で選別すると，CIM の生成を含めた同じ総時間で到達品質が上がるか．

本論文の貢献は次のとおりである．
\begin{itemize}
  \item 同じカット値（と 1 反転局所探索の改善量）の CIM 解の間に，後段 SA での価値の再現性のある差があることを，G-set\cite{helmberg2000}で示した．差の有無は SA の設定に強く依存し，SA が種から大きく離れる設定（開始温度が高い，予算が長い）では差が消える（\ref{sec:prereg}, \ref{sec:gset} 節）．
  \item 探索的な解析で見つけた固有ベクトル指標を事前登録して新しいデータで検証し，G22 では再現することを確かめた．一方，小規模問題の検証用データと他の G-set 問題では効き方が一様でなく，局所的な指標や短い SA の試し打ちのほうが安定していた（\ref{sec:prereg}--\ref{sec:gset} 節）．
  \item 18〜24 頂点の 30 問題の全列挙により，種の価値は最適解までの距離（経路長）でほぼ決まり（層内 Kendall $\tau\approx-0.74$），エネルギー障壁はそれより弱いことを示した（\ref{sec:small} 節）．
  \item CIM の生成，指標の計算，後段の探索をすべて含めた同じ総時間での比較を行い，選別の実用上の効果を評価した（\ref{sec:selection} 節）．
  \item 付随して，メメティックアルゴリズムのタブー探索で同点の手を常に番号最小の頂点から選ぶ実装が，疎なグラフで探索を止めることを見つけ，修正と高速化を行った（\ref{sec:ga} 節）．
\end{itemize}

研究の進め方として，探索的に得た結果は，判定基準と解析を固定した事前登録の文書をリポジトリにコミットしてから新しいデータで確かめた．登録どおりの結果は，仮説を支持しないものも含めてすべて報告する．
